\subsection{Numerical incompleteness}
\label{sec:numerical-incompleteness}
\label{subsec:numerical-incompleteness}

Use the numerical realization and comparison functors of
Section~\ref{subsec:numerical-realization}. Let
$I:\mathrm{Rec}\to\mathscr S_{\mathrm{rec}}$ be the full restriction
functor $I(X)=(X,1_X)$, $I(f)=f$. It preserves the restriction
idempotents on each object of $\mathrm{Rec}$, their order, and their
complements. Put $E=I(\omega)=Q(V)$.

The denotation
$\tau_{\mathrm{rec}}=\varphi_u:\omega\times\omega\rightharpoonup\omega$
is a Turing morphism in $\mathrm{Rec}$ by the specialization identity
of Lemma~\ref{lem:arithmetized-evaluation}. Let the isomorphism
$\zeta:\omega\to\omega+\omega$ in $\mathrm{Rec}$ send $2m$ to
the image of $m$ under the first coproduct injection and $2m+1$
to the image of $m$ under the second. The formulas $\alpha,\beta$ of
Section~\ref{subsec:arithmetic-internal-comparison} define the sets
$D_0,D_1$ of inputs on which the diagonal computation has even and
odd output, respectively. The partial identities
$1_{D_0},1_{D_1}:\omega\rightharpoonup\omega$ in $\mathrm{Rec}$ are
therefore the two diagonal restriction idempotents for
$\tau_{\mathrm{rec}}$ and $\zeta$.

The functor $J_{\mathrm{rec}}:\mathrm{Rec}\to\Pfn$ has
$J_{\mathrm{rec}}(X)=|X|$ and $J_{\mathrm{rec}}(f)=f$.
It preserves restrictions, ranges, zero morphisms, and the chosen
products and coproducts by their numerical descriptions.
In $\Pfn$, inverse images along total morphisms preserve binary joins of restriction
idempotents.
For complementary restriction idempotents $c,c':W\rightharpoonup W$
in $\mathrm{Rec}$, let $C,C'\subseteq|W|$ be their domains. For
$x\in|W|$, parallel enumeration of $C$ and $C'$ in the coding
$\gamma_W:|W|\to\N$ decides which contains $x$. The morphism
$s:W\to W+W$ in $\mathrm{Rec}$
given by
\[
 s(x)=
 \begin{cases}
  i_1(x),&x\in C,\\
  i_0(x),&x\in C'
 \end{cases}
\]
is total recursive; here $i_0,i_1:W\to W+W$ are the coproduct
injections. Its restrictions to the two domains give
$sc=i_1c$ and $sc'=i_0c'$. Since $J_{\mathrm{rec}}(1_\omega)\ne0$,
Corollary~\ref{cor:source-conditionals} gives inseparability
of $1_{D_0},1_{D_1}$ in $\mathrm{Rec}$. A recursive separator of $D_0,D_1$ would
give complementary restriction idempotents separating $1_{D_0},1_{D_1}$, so
$D_0,D_1$ are recursively inseparable. A restriction idempotent on
$E=(\omega,1_\omega)$ in $\mathscr S_{\mathrm{rec}}$ is precisely
a partial identity on a c.e.\ subset of $\N$, with meet and join
given by intersection and union. Hence any complementary pair
separating $I(1_{D_0}),I(1_{D_1})$ would separate $1_{D_0},1_{D_1}$ in $\mathrm{Rec}$.

For a c.e.\ arithmetic presentation $T$ containing the PA sentence
axioms, use $S_{\mathrm{rec}},\nu,\theta,t_T$ from
Section~\ref{subsec:numerical-realization}. The code comparison there
identifies their underlying numerical functions and domains with
the standard interpretations in
Lemma~\ref{lem:arithmetic-diagonal-representation}. For $n\in D_0$, PA proves $\alpha(\bar n)$, and for $n\in D_1$, it
proves $\neg\alpha(\bar n)$. Since $T$ contains the PA axioms, we obtain
\begin{equation}\label{eq:arithmetic-domain-inclusions}
 I(1_{D_0})\le\theta^*t_{\PA},\qquad
 I(1_{D_1})\le\theta^*\nu^*t_{\PA},\qquad t_{\PA}\le t_T.
\end{equation}
In $\mathscr S_{\mathrm{rec}}$, the first two inequalities compare restriction
idempotents on $E$ and the third compares restriction idempotents on
$S_{\mathrm{rec}}$.
For consistent $T$, its sentence theorem set and refutation set
are disjoint, so $t_T\nu^*t_T=0_{S_{\mathrm{rec}}}$.
Joins in this splitting are partial identities on unions of c.e.\ subsets,
and inverse images along total morphisms preserve these joins. Corollary~\ref{cor:domain-incompleteness}
therefore gives inseparability of $t_T,\nu^*t_T$ and
$t_T\vee\nu^*t_T\ne1_{S_{\mathrm{rec}}}$.

